% Appendix sub-section
\subsection{Implementation Details} \label{implementation}
\subsubsection{Wikipedia clone detection} \label{clones}
For detecting whether a reference document, $d$, is a Wikipedia article clone
we compute the
maximum recall of unigrams for each section of the Wikipedia article, $a$:

\[
r(d, a) = \max_{s \in sections(a)} \frac{|unigrams(d) \cap unigrams(s)|}{|unigrams(s)|} \]

and detect a clone if $r>0.5$.


\subsection{Human Evaluation experiment} \label{human_eval}
\begin{figure}[h]
\begin{center}
\fbox{%
    \includegraphics[width=0.7\textwidth]{ling_qual.png}
}%
\caption{Screenshot of DUC-style linguistic quality human evaluation tool.}
\label{crowdcompute_ling}
\end{center}
\end{figure}
To assess linguistic quality, we randomly selected samples
generated by models from the test set and ask raters to choose a score 
from 1 to 5 (higher is better) for five dimensions: Grammaticality,
Non-redundancy, Referential clarity, Focus, and Structure and Coherence. These
were used in the past at DUC for evaluating summaries \citep{dang2005overview}. For each model we selected 25 examples and averaged the scores for each
question across 3 raters (out of pool of 7).


\begin{figure}[h]
\begin{center}
\fbox{%
    \includegraphics[width=0.7\textwidth]{crowdcompute.png}
}%
\caption{Screenshot of side-by-side human evaluation tool. Raters are asked whether they prefer
model output on the left or right, given a ground truth Wikipedia text.}
\label{crowdcompute}
\end{center}
\end{figure}
To compare two models by human evaluation, we randomly select examples from the test set and
show model outputs side-by-side in the interface shown in Figure \ref{crowdcompute}. 
Which side a model appears on is randomized per example and rater. For the
experiments in Table \ref{human_eval_table} we had 3 raters score 25 examples each
and computed the ratio of ratings preferring one model over the other.

